% !TeX root = ../papers/nc-p3.tex
\section{単純表現の周期尾部と三行列恒等式の全正則窓}
\label{sec:periodic-tail-vertex-bridges}
\begin{lemma}\label{lem:periodic-tail-vertex-bridges}
有限四辺の標準正則窓で、頂点と両向きに全Extが零な均衡単純 $D$ を取り、
その各頂点次元を $c$ とする。均衡単純 $S$ の各頂点次元を $d$ とし、
$\operatorname{Ext}^1(k_i,S)=0$ が全ての $i$ で成立する一方、
$a_i=\dim\operatorname{Ext}^1(S,k_i)$ の少なくとも一つが正とする。
任意の頂点 $e$ の隅は $M_{c+d}(k)$ の全恒等式を満たすことができない。
全ての向きを逆にしても同じである。
\end{lemma}
\begin{proof}
$P_j$ を $D$ の周期表現を整数次数 $j$ から右へ展開した右加群とする。
次数 $n\ge j$ の空間は $D_{n\bmod4}$ で、矢は $D$ の作用である。
次数を忘れると $P_{j+4}\simeq P_j$、尾部の包含を一周したものを $z:P_j\to P_j$ と書く。
$P_j/(z-\lambda)P_j$ は $\lambda\ne0$ で $D$ の一周倍率を変えたもの、$\lambda=0$ で
頂点単純だけを持つ四層の加群である。

$P_j$ は有限生成射影による長さ3以下の分解を持つ。左頂点のKoszul分解に $P_j$ を
掛けて終端次数ごとに見ると、$n\ge j+4$ では全項が揃い、$D$ とその終端の頂点単純との
Tor複体そのものになる。これは仮定で完全である。残る有限次数には第四項がない。
下に有界な次数付き加群の極小分解を作れば主張を得る。
従って $E_j^m=\operatorname{Ext}^m(P_j,S)$ は次数を忘れても各々有限次元である。
$E_j^0=E_j^3=0$ であることを確かめる。$P_j/(z-\lambda)P_j$ を $D_\lambda$ と書くと、
全ての $\lambda$ で $\operatorname{Hom}(D_\lambda,S)=\operatorname{Hom}(S,D_\lambda)=0$。
非零倍率では頂点とのExt消滅が $D_\lambda$ と $S$ を区別し、零倍率では組成因子が異なるからである。
$z-\lambda$ の単射の完全列と四次元CY双対により、$E_j^0$ 上の $z-\lambda$ は単射、
$E_j^3$ 上では全射である。代数閉体で全ての $\lambda$ にこれが成り立つため両空間は零となる。

$\chi(P_j,S)=0$ である。実際有限極小分解の各頂点射影から $S$ へのHom次元は一律に $d$。
符号付き生成元数は $P_j$ のHilbert級数に $(1-t)^4$ を掛けて $t=1$ とした値で、
$ct^j(1-t)^3$ はそこで消える。よって $\dim E_j^1=\dim E_j^2$。
$0\to P_{j+1}\to P_j\to k_j^c\to0$ より $E_j^1\to E_{j+1}^1$ は単射で、
四周期性と有限次元性から全て同型である。続く完全列で
$\operatorname{Ext}^2(k_j^c,S)$ が $E_j^2$ に単射で入り、これが次の制限写像の核である。
その次元は $c a_j$ である（均衡性のEuler数零とCY双対）。
特に全ての $E_j^1,E_j^2$ は共通の正次元を持つ。
$z$ は $E_j^1$ 上可逆であり、ある $\lambda\ne0$ で $z-\lambda$ が核を持つ。
完全列から $\operatorname{Ext}^1(D_\lambda,S)\ne0$ を得る。

指定隅 $e$ に終わる部分も作る。ここでは別に、任意の点グラフの薄い単純を使って
上と同じ周期尾部 $P_j'$ を作る。薄い単純は全頂点と両向きの第一Extが消え、
均衡性とCY双対から全Extが消えるので、先の計算は $c=1$ として使える。
一周の $z$ は $\operatorname{Ext}^2(P_j',S)$ 上で非可逆である。
実際四つの制限の各項は同次元で、どの巡回順でもある $a_i>0$ の段に核がある。
$j=e+1\bmod4$ とし、$Q=P_j'/zP_j'$ を取る。$E_j^3=0$ と完全列により
$\operatorname{Ext}^3(Q,S)$ はこの非可逆写像の余核で、従って
$\operatorname{Ext}^1(S,Q)\ne0$ である。$Q$ の各頂点は一次元で末尾以外の矢は非零であり、
どの非零部分加群も末尾 $k_e$ を含む。

$D_\lambda$ と $Q$ の間の全Extは組成列から零である。従って上の二つの非分裂拡張の積は零で、
上から $D_\lambda,S,Q$ の三層加群に組める。有限像の半単純商の行列単位を頂点と両立して持ち上げる。
第一ブロックから $S$ への非零矩形があり、持ち上げた $S$ ブロックが下二層で生成する部分加群は
$S$ へ全射し、$Q$ に非零で交わる（交わらなければ分裂）。従って $k_e$ を含む。
$D_\lambda,S,k_e$ は互いに非同型かつ各一度であるため、中間の行列単位を挿入し、
$e$ 番における三ブロックを順に経る非零積を得る。
それらの次元和は $c+d+1$ であり、補題\ref{lem:pi-yoneda-bridges}の交代式の論証に反する。
\end{proof}

\begin{theorem}\label{thm:all-three-matrix-clifford}
有限四辺の標準正則窓の全四隅が四行列の恒等式を満たすとする。
三次元の均衡単純は頂点単純と両向きに全Extが零である。
どれか一隅が三行列の恒等式も満たせば、窓は必ず通常のClifford窓である。
従って全四隅に三行列の恒等式を課した正則窓の分類には半素性の仮定を全く要しない。
また、四行列の場合でも頂点と直交する非薄単純が一つあれば、全窓は左右Noetherで中心有限な合成整域となる。
特に三次元または四次元の均衡単純が一つあればこの結論が成立する。
\end{theorem}
\begin{proof}
全非頂点単純は均衡し、大きさは1から4である。四次元単純は隅の中心多項式で頂点から分離され、
薄い単純も既述の消滅を持つ。二・三次元では定理\ref{thm:one-sided-balanced-small-simples}で
拡張は高々片側に限られる。三次元で片側が残れば前補題を薄い $D$ に使い、$1+3+1>4$ で矛盾する。
二次元でも三行列の隅が指定されれば $1+2+1>3$ で矛盾する。
また直交する非薄な $D$ があれば全隅で $c+2+1>4$ となり同じ矛盾を得る。
これで全ての非頂点単純が頂点と両向きに直交し、定理\ref{thm:pi-vertex-extension-criterion}が
中心有限性、定理\ref{thm:noetherian-pi-window-cutoff}が合成整性を与える。
三行列の隅がある場合の指数は高々3であり、指数3は
系\ref{cor:no-index-three-finite-noetherian-window}で除かれ、有限辺では指数1もない。
指数2なら定理\ref{thm:one-corner-clifford-classification}から通常のClifford窓となる。
逆に同定理の基点無しで階数二以下との交叉が有限な系は全て対象に入る。
\end{proof}
